\setlength{\absparsep}{18pt}
\begin{resumo}
Este relatório parcial apresenta o andamento da ação de extensão que desenvolve um protótipo de bracelete inteligente baseado em microcontrolador ESP32, voltado a auxiliar pessoas com baixa visão na orientação em ambientes internos. Durante o detalhamento do projeto físico, a configuração do bracelete foi definida sem sensor de distância, e a detecção passou a ser feita por balizas \textit{Bluetooth Low Energy} (BLE) fixadas em pontos de interesse do ambiente, como portas e mesas. O bracelete, implementado em uma placa ESP32-C3, varre continuamente os anúncios das balizas, estima a distância de cada uma pela intensidade do sinal recebido (RSSI), aciona um motor de vibração conforme a baliza mais próxima e envia a lista de pontos detectados a um \textit{smartphone} em sentido único, por notificações BLE. No celular, uma página web que usa as interfaces Web Bluetooth e de síntese de voz apresenta os pontos em alto contraste e os anuncia em áudio, sem necessidade de instalar aplicativo. Até o momento foram concluídos o firmware do bracelete e da baliza, a interface web, um modo de demonstração e testes funcionais de comunicação. Os testes confirmaram a transmissão ESP32 $\rightarrow$ celular e a detecção simultânea de mais de uma baliza, mas mostraram que a estimativa de distância em metros pelo RSSI é instável com as placas utilizadas: com as placas a cerca de 1~m, a interface chegou a indicar 20~m. Como próxima etapa, a distância em metros será substituída por zonas de proximidade calibradas experimentalmente, seguidas da montagem física do bracelete e da validação com o público-alvo.

\textbf{Palavras-chave}: Tecnologia assistiva; ESP32; Bluetooth Low Energy; Baixa visão; Sistemas embarcados.
\end{resumo}

\begin{resumo}[Abstract]
\begin{otherlanguage*}{english}
This partial report presents the progress of an extension project that develops a smart wristband prototype based on the ESP32 microcontroller, intended to help people with low vision orient themselves indoors. While detailing the physical design, the wristband was defined without a distance sensor, so detection now relies on Bluetooth Low Energy (BLE) beacons placed at points of interest such as doors and tables. The wristband, built on an ESP32-C3 board, continuously scans beacon advertisements, estimates the distance to each one from the received signal strength (RSSI), drives a vibration motor according to the nearest beacon and sends the list of detected points to a smartphone in a single direction through BLE notifications. On the phone, a web page using the Web Bluetooth and speech synthesis interfaces shows the points in high contrast and announces them by voice, with no app installation. So far, the wristband and beacon firmware, the web interface, a demonstration mode and functional communication tests have been completed. The tests confirmed ESP32-to-phone transmission and simultaneous detection of more than one beacon, but showed that RSSI-based distance in meters is unstable with the boards used: with the boards about 1~m apart, the interface reported up to 20~m. As next steps, distance in meters will be replaced by experimentally calibrated proximity zones, followed by the physical assembly of the wristband and validation with the target audience.

\vspace{\onelineskip}
\noindent
\textbf{Keywords}: Assistive technology; ESP32; Bluetooth Low Energy; Low vision; Embedded systems.
\end{otherlanguage*}
\end{resumo}
